\section{Resumen y entorno de medición}

Se implementaron tres versiones en C (C17) con OpenMP de una simulación gravitacional N-Body en dos
dimensiones y se analizaron con la metodología PCAM:
\begin{itemize}
  \item \textbf{A} --- \code{nbody\_seq.c}: versión secuencial (línea base).
  \item \textbf{B} --- \code{nbody\_parallel.c}: paralela directa; cada cuerpo acumula su propia fuerza.
  \item \textbf{C} --- \code{nbody\_optimized.c}: cada par $(i,j)$ se calcula una sola vez usando
        $F_{ij} = -F_{ji}$.
\end{itemize}

Todas las versiones usan \code{seed = 42}, $N = 5000$ cuerpos, 10 pasos y $\Delta t = 0.01$, y miden el
tiempo con \code{omp\_get\_wtime()}. Cada configuración se ejecutó tres veces y se reporta el promedio
(\code{results.csv}, generado por \code{run\_experiments.sh}). El código fuente completo está en el repositorio \repo.

\begin{table}[H]
\centering
\caption{Entorno de medición.}
\begin{tabular}{@{}ll@{}}
\toprule
Procesador     & Apple M2 Pro, 10 núcleos (6 de rendimiento + 4 de eficiencia) \\
Memoria        & 16 GB \\
OpenMP         & \code{omp\_get\_num\_procs()} $= 10$ \\
Compilador     & GCC 14.2, \code{-O2 -std=c17 -fopenmp} \\
\bottomrule
\end{tabular}
\end{table}

Los 8 hilos se reparten entre núcleos heterogéneos y una corrida completa dura menos de 0.4\,s, por lo
que la variación entre corridas es de varios por ciento; diferencias menores a $\sim$5\,\% no se
consideran concluyentes.

La mejor configuración obtenida fue la versión C con \code{reduction}, \code{schedule(dynamic, 64)} y
8 hilos: 0.069\,s frente a 0.355\,s de la secuencial, es decir, $S = 5.17$ y $E = 0.65$, con el mismo
checksum que la versión secuencial.
